\slcol{सञ्जय उवाच —\\
\Index{तं तथा कृपया}विष्टमश्रुपूर्णाकुलेक्षणम्।\\
 विषीदन्तमिदं वाक्यमुवाच मधुसूदनः॥ २.१॥}
\translate{sañjaya uvāca —\\
taṁ tathā kṛpayāviṣṭamaśrupūrṇākulēkṣaṇam |\\
viṣīdantamidaṁ vākyamuvāca madhusūdanaḥ ‖ 2.1 ‖}
\cquote{Sanjaya said: Seeing Arjuna overwhelmed by pity and despondency, his eyes filled with tears and mind bewildered, Madhusudana (Krishna) spoke the following words to him.}
\slcol{श्रीभगवानुवाच —\\
\Index{कुतस्त्वा कश्मलमिदं} विषमे समुपस्थितम्।\\
 अनार्यजुष्टमस्वर्ग्यमकीर्तिकरमर्जुन॥ २.२॥}
\translate{śrībhagavānuvāca —\\
kutastvā kaśmalamidaṁ viṣamē samupasthitam |\\
anāryajuṣṭamasvargyamakīrtikaramarjuna ‖ 2.2 ‖}
\cquote{Sri Bhagavan said: Whence has this weakness come upon you at this critical moment, O Arjuna? It is unworthy of one noble (\emph{arya}), and leads not to heaven—this faint-heartedness brings only disgrace.}

\newpage
\slcol{\Index{क्लैब्यं मा स्म गमः} पार्थ नैतत्त्वय्युपपद्यते।\\
क्षुद्रं हृदयदौर्बल्यं त्यक्त्वोत्तिष्ठ परन्तप॥ २.३॥}
\translate{klaibyaṁ mā sma gamaḥ pārtha naitattvayyupapadyatē |\\
kṣudraṁ hṛdayadaurbalyaṁ tyaktvōttiṣṭha parantapa ‖ 2.3 ‖}
\cquote{Yield not to this unmanliness, O Pārtha (Arjuna)! It does not befit you. Cast off this petty weakness of heart, O Pārantapa (scorcher of foes), and arise!}
\slcol{अर्जुन उवाच —\\
\Index{कथं भीष्ममहं} सं‍ख्ये द्रोणं च मधुसूदन।\\
इषुभिः प्रतियोत्स्यामि पूजार्हावरिसूदन॥ २.४॥}
\translate{arjuna uvāca —\\
kathaṁ bhīṣmamahaṁ saṁ‍khyē drōṇaṁ ca madhusūdana |\\
iṣubhiḥ pratiyōtsyāmi pūjārhāvarisūdana ‖ 2.4 ‖}
\cquote{Arjuna said: O Mādhusūdana, how can I face Bhīṣma and Droṇa with arrows in battle—they who are worthy of reverence, O Arisudana (destroyer of foes)?}
\slcol{\Index{गुरूनहत्वा हि} महानुभावान् श्रेयो भोक्तुं भैक्षमपीह लोके।\\
 हत्वार्थकामांस्तु गुरूनिहैव भुञ्जीय भोगान्रुधिरप्रदिग्धान्॥ २.५॥}
\translate{gurūnahatvā hi mahānubhāvān śrēyō \\
\hspace*{0.5cm}bhōktuṁ bhaikṣamapīha lōkē |\\
hatvārthakāmāṁstu gurūnihaiva \\
\hspace*{0.5cm}bhuñjīya bhōgānrudhirapradigdhān ‖ 2.5 ‖}
\cquote{Better it is, I feel, to live on alms in this world than to slay these noble teachers. By killing them, even though they are my enemies, I would be enjoying pleasures and wealth stained with blood.}
\slcol{\Index{न चैतद्विद्मः कतरन्नो} गरीयो यद्वा जयेम यदि वा नो जयेयुः।\\
यानेव हत्वा न जिजीविषामस्तेऽवस्थिताः प्रमुखे धार्तराष्ट्राः॥ २.६॥}
\translate{na caitadvidmaḥ katarannō garīyō\\
\hspace*{0.5cm}yadvā jayēma yadi vā nō jayēyuḥ |\\
yānēva hatvā na jijīviṣāmastē\\
\hspace*{0.5cm}vasthitāḥ pramukhē dhārtarāṣṭrāḥ ‖ 2.6 ‖}

\newpage
\begin{mananam}{\mananamfont {Mananam sloka - 2, 3}}
\mananamtext Do I tend to wallow in self-pity? Do certain outer conditions or inner urges pull down my energies, and do I yield to my weaknesses? Am I too sensitive to shut out corrections by my teachers or well-meaning friends?
Do I give up easily on my goals and resolutions? What does it take for me to bring about the intended changes in my life?
\end{mananam}
\WritingHand\enspace{\selfReflect\small{Self Reflection}}
\begin{inspiration}{\mananamfont Inspiration}
\mananamtext You are never too weak to face your challenges. Don’t give into doubts, self-pity and temptations. Get up and get going with your highest goals and aspirations. When you make a firm determination, all the positive forces of the universe will support and strengthen you.
\end{inspiration}
\newpage

\cquote{We do not know which is better for us, whether we should conquer them, or they should conquer us. Those very sons of Dhṛitarāṣṭra stand before us, after whose death we would not wish to live.}
\slcol{\Index{कार्पण्यदोषोपहतस्वभावः} \\
\hspace*{1.0cm} पृच्छामि त्वां धर्मसंमूढचेताः।\\
यच्छ्रेयः स्यान्निश्चितं ब्रूहि तन्मे शिष्यस्तेऽहं \\
\hspace*{1.0cm} शाधि मां त्वां प्रपन्नम्॥ २.७॥}
\translate{kārpaṇyadōṣōpahatasvabhāvaḥ\\
\hspace*{0.5cm} pṛcchāmi tvāṁ dharmasaṁmūḍhacētāḥ |\\
yacchrēyaḥ syānniścitaṁ brūhi tanmē\\
\hspace*{0.5cm} śiṣyastēhaṁ śādhi māṁ tvāṁ prapannam ‖ 2.7 ‖}
\cquote{My very nature is stricken with weakness; my mind is utterly confused about \emph{dharma}. I ask You to tell me what is truly good (\emph{shreyas}) for me. I am Your disciple; I take refuge in You; please instruct me.}
\slcol{\Index{न हि प्रपश्यामि} ममापनुद्याद्यच्छोकमुच्छोषणमिन्द्रियाणाम्।\\
अवाप्य भूमावसपत्नमृद्धं राज्यं सुराणामपि चाधिपत्यम्॥ २.८॥}
\translate{na hi prapaśyāmi\\ 
\hspace*{0.5cm}mamāpanudyādyacchōkamucchōṣaṇamindriyāṇām |\\
avāpya bhūmāvasapatnamṛddhaṁ rājyaṁ\\
\hspace*{0.5cm} surāṇāmapi cādhipatyam ‖ 2.8 ‖}
\cquote{Even if I were to obtain a prosperous and unrivalled kingdom on earth, or lordship over the gods in heaven, I do not see what could dispel this sorrow that dries up my senses.}
\slcol{सञ्जय उवाच —\\
\Index{एवमुक्त्वा हृषीकेशं} गुडाकेशः परन्तपः।\\
न योत्स्य इति गोविन्दमुक्त्वा तूष्णीं बभूव ह॥ २.९॥}
\translate{sañjaya uvāca —\\
ēvamuktvā hṛṣīkēśaṁ guḍākēśaḥ parantapaḥ |\\
na yōtsya iti gōvindamuktvā tūṣṇīṁ babhūva ha ‖ 2.9 ‖}
\cquote{Sanjaya said: Having spoken thus to Hṛṣīkeśa (Krishna), the destroyer of foes, Guḍākeśa (Arjuna) said, “I shall not fight,” and became silent.}

\newpage
\begin{mananam}{\mananamfont Mananam sloka - 7}
\mananamtext With right understanding, awareness of our weakness and limitations can be transformed into spiritual traits of humility and reverence. 
Am I humble enough to accept that I don’t know everything? How can I open my mind and be receptive to the life-transforming truths of the rishis and the spiritually evolved guides in my life?
\end{mananam}
\WritingHand\enspace{\selfReflect\small{Self Reflection}}
\begin{inspiration}{\mananamfont Inspiration}
\mananamtext “When the student is ready, the teacher will come” is an adage told to all spiritual seekers. The readiness is not one based in time but rather is found in one’s mental states of open mindedness and receptivity.
True humility is not a sign of weakness, but of greatness.
\end{inspiration}
\newpage

\slcol{\Index{तमुवाच हृषीकेशः} प्रहसन्निव भारत।\\
सेनयोरुभयोर्मध्ये विषीदन्तमिदं वचः॥ २.१०॥}
\translate{tamuvāca hṛṣīkēśaḥ prahasanniva bhārata |\\
sēnayōrubhayōrmadhyē viṣīdantamidaṁ vacaḥ ‖ 2.10‖}
\cquote{Then Hṛṣīkeśa, smiling, spoke the following words to the despondent Arjuna, grieving amidst both armies.}
\slcol{श्रीभगवानुवाच —\\
\Index{अशोच्यानन्वशोचस्त्वं} प्रज्ञावादांश्च भाषसे।\\
 गतासूनगतासूंश्च नानुशोचन्ति पण्डिताः॥ २.११॥}
\translate{śrībhagavānuvāca —\\
aśōcyānanvaśōcastvaṁ prajñāvādāṁśca bhāṣasē |\\
gatāsūnagatāsūṁśca nānuśōcanti paṇḍitāḥ ‖ 2.11‖}
\cquote{Sri Bhagavan said: You grieve for those who should not be grieved for, yet you speak words of wisdom. The wise do not mourn for the living or the dead.}
\slcol{\Index{न त्वेवाहं जातु} नासं न त्वं नेमे जनाधिपाः।\\
न चैव न भविष्यामः सर्वे वयमतः परम्॥ २.१२॥}
\translate{na tvēvāhaṁ jātu nāsaṁ na tvaṁ nēmē janādhipāḥ |\\
na caiva na bhaviṣyāmaḥ sarvē vayamataḥ param ‖ 2.12‖}
\cquote{Never was there a time indeed when I did not exist, nor you, nor these kings; nor shall any of us cease to be hereafter.}
\slcol{\Index{देहिनोऽस्मिन्यथा} देहे कौमारं यौवनं जरा।\\
तथा देहान्तरप्राप्तिर्धीरस्तत्र न मुह्यति॥ २.१३॥}
\translate{dēhinō'sminyathā dēhē kaumāraṁ yauvanaṁ jarā |\\
tathā dēhāntaraprāptirdhīrastatra na muhyati ‖ 2.13‖}
\cquote{Just as the embodied soul passes through childhood, youth, and old age in this body, so too it passes into another body after death. The wise are not deluded by this.}

\newpage
\begin{mananam}{\mananamfont Mananam sloka - 11}
\mananamtext Do I tend to use authoritative teachings to endorse my stance? In difficult situations of life, am I brave enough to do that which is difficult, or do I tend to take the easy way out?
How can I keep alive the spirit of those who pass away? Can I learn to see their virtues and to honour their spirit by practicing those values in my own life?
\end{mananam}
\WritingHand\enspace{\selfReflect\small{Self Reflection}}
\begin{inspiration}{\mananamfont Inspiration}
\mananamtext To walk the talk is a sign of mental strength. To give excuses for one’s inactions by quoting scriptures or the examples of others is a sign of weakness.
\end{inspiration}
\newpage

\newpage
\begin{mananam}{\mananamfont Mananam sloka - 13}
\mananamtext Do I resist changes in my life? Does youthful energy and good health make me elated, whereas a drop in energy and the onset of ill health upset me? What does it mean to me to remain mentally steady amidst all physical transformations, knowing that everything in nature is constantly undergoing changes? Is it possible that I self-impose on myself limitations in life based on my physical features and age?
\end{mananam}
\WritingHand\enspace{\selfReflect\small{Self Reflection}}
\begin{inspiration}{\mananamfont Inspiration}
\mananamtext The natural process of this body is to go through childhood, youth and old age. A sure means of sorrow is to associate oneself with only the beauty of youth. Far superior to physical beauty is mental beauty (purity). Superior even to that is the beauty of \emph{atman}, that changeless entity associated with the body.
In the Creator’s highest wisdom, our body has been designed to be changing so that we might realise our changeless nature to be beyond our physical entity.
\end{inspiration}
\newpage

\slcol{
\Index{मात्रास्पर्शास्तु कौन्तेय} शीतोष्णसुखदुःखदाः। \\
आगमापायिनोऽनित्यास्तांस्तितिक्षस्व भारत॥ २.१४॥}
\translate{mātrāsparśāstu kauntēya śītōṣṇasukhaduḥkhadāḥ |\\
āgamāpāyinō'nityāstāṁstitikṣasva bhārata ‖ 2.14‖}
\cquote{O son of Kuntī, the contacts of the senses with their objects produce cold and heat, pleasure and pain. They come and go; they are impermanent. Endure them bravely, O Bhārata (Arjuna).}
\slcol{\Index{यं हि न व्यथयन्त्येते} पुरुषं पुरुषर्षभ।\\
समदुःखसुखं धीरं सोऽमृतत्वाय कल्पते॥ २.१५॥}
\translate{yaṁ hi na vyathayantyētē puruṣaṁ puruṣarṣabha |\\
samaduḥkhasukhaṁ dhīraṁ sō'mṛtatvāya kalpatē ‖ 2.15‖}
\cquote{That person, O best among men (Arjuna), the wise one who is steady in pleasure and pain; whom these do not disturb, is fit for immortality.}
\slcol{\Index{नासतो विद्यते भावो} नाभावो विद्यते सतः।\\
उभयोरपि दृष्टोऽन्तस्त्वनयोस्तत्त्वदर्शिभिः॥ २.१६॥}
\translate{nāsatō vidyatē bhāvō nābhāvō vidyatē sataḥ |\\
ubhayōrapi dṛṣṭō'ntastvanayōstattvadarśibhiḥ ‖ 2.16‖}
\cquote{Of the unreal there is no existence; of the real there is no non-existence. The nature of both is perceived by the seers of truth.}
\slcol{\Index{अविनाशि तु तद्विद्धि} येन सर्वमिदं ततम्।\\
विनाशमव्ययस्यास्य न कश्चित्कर्तुमर्हति॥ २.१७॥}
\translate{avināśi tu tadviddhi yēna sarvamidaṁ tatam |\\
vināśamavyayasyāsya na kaścitkartumarhati‖ 2.17‖}
\cquote{Know that to be indestructible, by which all this is pervaded. None can bring about the destruction of this imperishable essence.}

\newpage
\begin{mananam}{\mananamfont Mananam sloka - 14}
\mananamtext Just as with our body, our mental states are also ever changing. Can we learn to be aware of our mental states without judging them as good or bad? Can we glimpse the subtle background of awareness which ever remains steady amidst our passing moods of anger, fear, sadness, excitement, contentedness and happiness?
Do I ever feel that I can no longer endure something? Do I become easily impatient and intolerant? If so, then it is most likely that I have taken the passing things in life to be permanent.
\end{mananam}
\WritingHand\enspace{\selfReflect\small{Self Reflection}}
\begin{inspiration}{\mananamfont Inspiration}
\mananamtext In the grand scheme of things, we weren’t designed to dwell on past memories of this life, but to dwell on the lost memory of our innate soul nature. Don’t waste time begging from life the petty experiences of fleeting happiness. Everlasting joy is your birthright!
\end{inspiration}
\newpage

\newpage
\begin{mananam}{\mananamfont Mananam sloka - 15}
\mananamtext Where do I experience happiness? Where do I experience sadness? Can I see through my own direct insight that these are all just mental states? Observe in your life those situations in which one object can bring sorrow to you but happiness to someone else, whether it is food, weather or anything else. Furthermore, can you observe the very awareness of your mental state itself? Does it come and go?
\end{mananam}
\WritingHand\enspace{\selfReflect\small{Self Reflection}}
\begin{inspiration}{\mananamfont Inspiration}
\mananamtext Immortality is the secret desire of every living being, just as death is the greatest fear. One who realizes the Self to be beyond body and mind is not troubled by his or her changing condition. That person becomes immortal; realizes that he or she was in essence ever immortal!
\end{inspiration}
\newpage

\newpage
\begin{mananam}{\mananamfont Mananam sloka - 16}
\mananamtext Am I brave enough to enquire into what is real and what is unreal in every situation of life? Is my perception of people and situations shaped by my past experiences, society’s beliefs, cultural norms, etc.? Am I willing to question my beliefs and limited logical reasonings? Am I willing to consider the truths extolled by sages and discover for myself the higher reality that brings ultimate freedom? 
Am I willing to commit my life to the quest for truth and reality?
\end{mananam}
\WritingHand\enspace{\selfReflect\small{Self Reflection}}
\begin{inspiration}{\mananamfont Inspiration}
\mananamtext There are various levels of relative truths. But there is only one ultimate truth. Only those who find the ultimate truth are free. All others are seeking that freedom, without realizing its nature. Once you commit your life to the quest for truth, then you gain great mental strength and power, just like King Harishchandra and Lord Rama.
\end{inspiration}
\newpage

\slcol{\Index{अन्तवन्त इमे देहा} नित्यस्योक्ताः शरीरिणः।\\
अनाशिनोऽप्रमेयस्य तस्माद्युध्यस्व भारत॥ २.१८॥}
\translate{antavanta imē dēhā nityasyōktāḥ śarīriṇaḥ |\\
anāśinō'pramēyasya tasmādyudhyasva bhārata ‖ 2.18‖}
\cquote{The bodies of the embodied soul are said to have an end. But the indwelling Self is eternal, indestructible, and unknowable. Therefore, fight, O Bhārata (Arjuna).}
\slcol{\Index{य एनं वेत्ति हन्तारं} यश्चैनं मन्यते हतम्।\\
उभौ तौ न विजानीतो नायं हन्ति न हन्यते॥ २.१९॥}
\translate{ya ēnaṁ vētti hantāraṁ yaścainaṁ manyatē hatam |\\
ubhau tau na vijānītō nāyaṁ hanti na hanyatē ‖ 2.19‖}
\cquote{He who thinks the Self kills, and he who thinks the Self is killed—both are ignorant. The Self neither kills nor is killed.}
\slcol{\Index{न जायते म्रियते} वा कदाचिन्नायं भूत्वाभविता वा न भूयः।\\
अजो नित्यः शाश्वतोऽयं पुराणो न हन्यते हन्यमाने शरीरे॥ २.२०॥}
\translate{na jāyatē mriyatē vā kadācinnāyaṁ\\
\hspace*{0.5cm}bhūtvābhavitā vā na bhūyaḥ |\\
ajō nityaḥ śāśvatō'yaṁ purāṇō\\
\hspace*{0.5cm}na hanyatē hanyamānē śarīrē ‖ 2.20‖}
\cquote{The Self is never born, nor does it ever die. It has never come into being, nor will it ever cease to be. Unborn, eternal, everlasting, and primeval, it is not slain when the body is slain.}
\slcol{\Index{वेदाविनाशिनं नित्यं} य एनमजमव्ययम्।\\
कथं स पुरुषः पार्थ कं घातयति हन्ति कम्॥ २.२१॥}
\translate{vēdāvināśinaṁ nityaṁ ya ēnamajamavyayam |\\
kathaṁ sa puruṣaḥ pārtha kaṁ ghātayati hanti kam ‖ 2.21‖}
\cquote{He who knows the Self to be indestructible, unborn, and eternal—how can that person think, “I kill,” or “I am killed”?}

\newpage
\begin{mananam}{\mananamfont Mananam sloka - {20, 21}}
\mananamtext Is my notion of myself limited to this physical body? Do I limit my perception of others to their physical existence alone? Can I try to tune into the thoughts and ever-living presence of great ones who are no longer alive? Can I intuit the changeless essence of myself and others, beyond their changing body and mind?
\end{mananam}
\WritingHand\enspace{\selfReflect\small{Self Reflection}}
\begin{inspiration}{\mananamfont Inspiration}
\mananamtext The feeling that we were born on such and such day and will die one day is a shared belief. When our intuition and understanding gets refined, we wake up to the reality that our essence, the soul, is not only beyond death, but was never born!
\end{inspiration}
\newpage

\slcol{\Index{वासांसि जीर्णानि} यथा विहाय नवानि गृह्णाति नरोऽपराणि।\\
तथा शरीराणि विहाय जीर्णान्यन्यानि संयाति नवानि देही॥ २.२२॥}
\translate{vāsāṁsi jīrṇāni yathā vihāya navāni gṛhṇāti narō'parāṇi | \\
tathā śarīrāṇi vihāya jīrṇānyanyāni saṁyāti navāni dēhī ‖ 2.22‖}
\cquote{Just as a person casts off worn-out garments and puts on new ones, so the embodied Self casts off worn-out bodies and enters into new ones.}
\slcol{\Index{नैनं छिन्दन्ति} शस्त्राणि नैनं दहति पावकः।\\
न चैनं क्लेदयन्त्यापो न शोषयति मारुतः॥ २.२३॥}
\translate{nainaṁ chindanti śastrāṇi nainaṁ dahati pāvakaḥ |\\
na cainaṁ klēdayantyāpō na śōṣayati mārutaḥ ‖ 2.23‖}
\cquote{Weapons cannot cut it, fire cannot burn it, water cannot wet it, nor can the wind dry it.}
\slcol{\Index{अच्छेद्योऽयमदाह्योऽयम}क्लेद्योऽशोष्य एव च।\\
नित्यः सर्वगतः स्थाणुरचलोऽयं सनातनः॥ २.२४॥}
\translate{acchēdyō'yamadāhyō'yamaklēdyō'śōṣya ēva ca |\\
nityaḥ sarvagataḥ sthāṇuracalō'yaṁ sanātanaḥ ‖ 2.24‖}
\cquote{This Self is said to be uncuttable, unburnable, impervious to water, and undryable. It is eternal, all-pervading, stable, immovable, and everlasting.}
\slcol{\Index{अव्यक्तोऽयमचिन्त्योऽयम}विकार्योऽयमुच्यते।\\
तस्मादेवं विदित्वैनं नानुशोचितुमर्हसि॥ २.२५॥}
\translate{avyaktō'yamacintyō'yamavikāryō'yamucyatē |\\
tasmādēvaṁ viditvainaṁ nānuśōcitumarhasi ‖ 2.25‖}
\cquote{This Self is said to be unmanifest, inconceivable, and immutable. Knowing this, you should not grieve for the body.}
\slcol{\Index{अथ चैनं नित्यजातं} नित्यं वा मन्यसे मृतम्।\\
तथापि त्वं महाबाहो नैवं शोचितुमर्हसि॥ २.२६॥}
\translate{atha cainaṁ nityajātaṁ nityaṁ vā manyasē mṛtam |\\
tathāpi tvaṁ mahābāhō naivaṁ śōcitumarhasi ‖ 2.26‖}


\newpage
\begin{mananam}{\mananamfont Mananam sloka - 22}
\mananamtext In our society, people are judged based on their outer appearance. Regardless of how many different clothes an actor might adorn based on the character that he projects, he remains, nevertheless, the person that he is. Can I learn to see people for what they are? Can I learn to see the pure soul-essence in all, despite their flaws and weaknesses?
\end{mananam}
\WritingHand\enspace{\selfReflect\small{Self Reflection}}
\begin{inspiration}{\mananamfont Inspiration}
\mananamtext A person of subtle vision learns to see other people not based on their outer appearance but based on their character. A saint sees everyone to be of their own soul image.
\end{inspiration}
\newpage

%% \newgeometry{margin=0pt} % Apply margin only for this page
%% \thispagestyle{empty}
%% \begin{figure}
%% \centering
%% \includegraphics[width=\paperwidth, height=\paperheight, keepaspectratio]{../images/002.jpg}
%% \end{figure}
%% \restoregeometry % Restore original geometry settings
%% \newpage

\newpage
\begin{mananam}{\mananamfont {Mananam sloka - 23, 24}}
\mananamtext As is my vision of myself, so is my attitude towards life and beyond. Am I afraid to face life’s circumstances, fearing challenges, difficulties, pain, and ultimately death? Or do I have an unshakeable vision of myself and others? Do I identify myself as being confined to the body?
\end{mananam}
\WritingHand\enspace{\selfReflect\small{Self Reflection}}
\begin{inspiration}{\mananamfont Inspiration}
\mananamtext Gross objects can affect gross elements but not subtle elements. Space is not affected by fire or by the air that passes through it. Soul is not affected by physical or mental transformations.
\end{inspiration}
\newpage

\cquote{Even if you think of the Self as constantly born and constantly dying, still, O mighty-armed (Arjuna), you should not grieve.}

\slcol{\Index{जातस्य हि ध्रुवो} मृत्युर्ध्रुवं जन्म मृतस्य च।\\
तस्मादपरिहार्येऽर्थे न त्वं शोचितुमर्हसि॥ २.२७॥}
\translate{jātasya hi dhruvō mṛtyurdhruvaṁ janma mṛtasya ca |\\
tasmādaparihāryē'rthē na tvaṁ śōcitumarhasi ‖ 2.27‖}
\cquote{For one who is born, death is certain; and for one who has died, birth is certain. Therefore, you should not grieve over the inevitable.}
\slcol{\Index{अव्यक्तादीनि भूतानि} व्यक्तमध्यानि भारत।\\
अव्यक्तनिधनान्येव तत्र का परिदेवना॥ २.२८॥}
\translate{avyaktādīni bhūtāni vyaktamadhyāni bhārata |\\
avyaktanidhanānyēva tatra kā paridēvanā ‖ 2.28‖}
\cquote{Beings are unmanifest in their beginning, manifest in their middle state, O Bharata (Arjuna), and unmanifest again at their end. Why then should one grieve over them?}
\slcol{\Index{आश्चर्यवत्पश्यति कश्चिदेन}माश्चर्यवद्वदति तथैव चान्यः।\\
आश्चर्यवच्चैनमन्यः शृणोति श्रुत्वाप्येनं वेद न चैव कश्चित्॥ २.२९॥}
\translate{āścaryavatpaśyati kaścidēnamāścaryavadvadati\\
\hspace*{0.5cm}tathaiva cānyaḥ |\\
āścaryavaccainamanyaḥ śṛṇōti śrutvāpyēnaṁ\\
\hspace*{0.5cm}vēda na caiva kaścit ‖ 2.29‖}
\cquote{Some behold It as a wonder; others speak of it as a wonder; others hear of It as a wonder; yet even after hearing, none truly understands It.}
\slcol{\Index{देही नित्यमवध्योऽयं} देहे सर्वस्य भारत।\\
तस्मात्सर्वाणि भूतानि न त्वं शोचितुमर्हसि॥ २.३०॥}
\translate{dēhī nityamavadhyō'yaṁ dēhē sarvasya bhārata |\\
tasmātsarvāṇi bhūtāni na tvaṁ śōcitumarhasi ‖ 2.30‖}
\cquote{The eternal indweller in every body is ever indestructible, O Bhārata. Therefore, you should not grieve for any creature.}

\newpage
\begin{mananam}{\mananamfont {Mananam sloka - 27}}
\mananamtext Do I fear ageing and growing old? Am I anxious about change? Can I make peace with changes in life? Can I bring in an accepting attitude toward the changing nature of my body and the environment around me?
\end{mananam}
\WritingHand\enspace{\selfReflect\small{Self Reflection}}
\begin{inspiration}{\mananamfont Inspiration}
\mananamtext “This too shall pass,” is an ancient adage. Whether in bad times or in good times, learning to reflect on this message will help instil in us a sense of acceptance and wellbeing.
\end{inspiration}
\newpage

\newpage
\begin{mananam}{\mananamfont {Mananam sloka - 29}}
\mananamtext Do I see and experience life with a fresh perspective? Am I biased, judgmental and critical of everything? Do I allow the experience of the present moment to be coloured by the past baggage?
\end{mananam}
\WritingHand\enspace{\selfReflect\small{Self Reflection}}
\begin{inspiration}{\mananamfont Inspiration}
\mananamtext Children look at everything with a sense of wonder and curiosity. There is a sense of freshness and joy in everything they experience. But as they grow up, they develop a bias towards what they want to hear and see. Soon life becomes dull and monotonous. Revive the child within you. Experience every moment of life with fresh eyes and an uncluttered mind.
\end{inspiration}
\newpage


\slcol{\Index{स्वधर्ममपि चावेक्ष्य} न विकम्पितुमर्हसि।\\
धर्म्याद्धि युद्धाच्छ्रेयोऽन्यत्क्षत्त्रियस्य न विद्यते॥ २.३१॥}
\translate{svadharmamapi cāvēkṣya na vikampitumarhasi |\\
dharmyāddhi yuddhācchrēyō'nyatkṣattriyasya na vidyatē ‖ 2.31‖}
\cquote{Considering one's own duty (\emph{swadharma}), you should not waver. For there is nothing higher for a \emph{kshatriya} (warrior) than a righteous battle.}
\slcol{\Index{यदृच्छया चोपपन्नं} स्वर्गद्वारमपावृतम्।\\
सुखिनः क्षत्रियाः पार्थ लभन्ते युद्धमीदृशम्॥ २.३२॥}
\translate{yadṛcchayā cōpapannaṁ svargadvāramapāvṛtam |\\
sukhinaḥ kṣatriyāḥ pārtha labhantē yuddhamīdṛśam ‖ 2.32‖}
\cquote{O son of Prithā (Arjuna), fortunate are the kshatriyas to whom such a righteous battle comes unsought, opening the very gates of heaven.}
\slcol{\Index{अथ चेत्त्वमिमं धर्म्यं} सङ्ग्रामं न करिष्यसि।\\
ततः स्वधर्मं कीर्तिं च हित्वा पापमवाप्स्यसि॥ २.३३॥}
\translate{atha cēttvamimaṁ dharmyaṁ saṅgrāmaṁ na kariṣyasi |\\
tataḥ svadharmaṁ kīrtiṁ ca hitvā pāpamavāpsyasi ‖ 2.33‖}
\cquote{But if you will not engage in this righteous battle, then by abandoning your own \emph{dharma} and honour, you will incur sin.}
\slcol{\Index{अकीर्तिं चापि भूतानि} कथयिष्यन्ति तेऽव्ययाम्।\\
सम्भावितस्य चाकीर्तिर्मरणादतिरिच्यते॥ २.३४॥}
\translate{akīrtiṁ cāpi bhūtāni kathayiṣyanti tē'vyayām |\\
sambhāvitasya cākīrtirmaraṇādatiricyatē ‖ 2.34‖}
\cquote{People will speak of your everlasting dishonour. For one who is well respected, dishonour is worse than death.}
\slcol{\Index{भयाद्रणादुपरतं मंस्यन्ते} त्वां महारथाः।\\
येषां च त्वं बहुमतो भूत्वा यास्यसि लाघवम्॥ २.३५॥}
\translate{bhayādraṇāduparataṁ maṁsyantē tvāṁ mahārathāḥ |\\
yēṣāṁ ca tvaṁ bahumatō bhūtvā yāsyasi lāghavam ‖ 2.35‖}


\newpage
\begin{mananam}{\mananamfont {Mananam sloka - 31}}
\mananamtext What are the duties in my life that I consider to be unfairly imposed on me? What are the brighter side of things if I learn to fulfil them? Isn’t it better that I learn to do all my life-given duties with an accepting attitude rather than with a complaining attitude?
\end{mananam}
\WritingHand\enspace{\selfReflect\small{Self Reflection}}
\begin{inspiration}{\mananamfont Inspiration}
\mananamtext A seemingly unpleasant duty equips us with some opportunity or skills. Such is the law of nature that nothing is ever unfair to anyone, and no endeavour ever goes unrewarded in the grand scale of time.
\end{inspiration}
\newpage


\newpage
\begin{mananam}{\mananamfont {Mananam sloka - 33}}
\mananamtext What is my \emph{dharma} in the present situation (based on my \emph{ashrama} and life responsibilities)? How do I embrace my \emph{dharma} with the changing times? To maintain self-respect and dignity in society, how should I behave? Am I abandoning any duties or values that I need to uphold? Do I fear facing my life’s battles and thereby put the fulfilment of my \emph{dharma} at risk?
\end{mananam}
\WritingHand\enspace{\selfReflect\small{Self Reflection}}
\begin{inspiration}{\mananamfont Inspiration}
\mananamtext To seek the supreme truth through renunciation is acceptable as long as one is wholeheartedly committed to it. To fall short in serving one’s family is acceptable if one is occupied with service of a more universal scope, e.g. serving the nation or a larger community. But to abandon one’s duties and responsibilities to family, society and to one’s own soul due to laziness and indulgence in one’s passions, is to be overcome with negativity and fears. Harbouring such wrong notions is the worst thing one can do to oneself and others.
\end{inspiration}
\newpage

\cquote{The great warriors who hold you in high esteem will think that you withdrew from battle out of fear. They will regard you lightly (look down on you).}

\slcol{\Index{अवाच्यवादांश्च} बहून्वदिष्यन्ति तवाहिताः।\\
निन्दन्तस्तव सामर्थ्यं ततो दुःखतरं नु किम्॥ २.३६॥}
\translate{avācyavādāṁśca bahūnvadiṣyanti tavāhitāḥ |\\
nindantastava sāmarthyaṁ tatō duḥkhataraṁ nu kim ‖ 2.36‖}
\cquote{Your enemies will speak of you with many unbefitting words, mocking your strength. What could be more painful than that?}
\slcol{\Index{हतो वा प्राप्स्यसि} स्वर्गं जित्वा वा भोक्ष्यसे महीम्।\\
तस्मादुत्तिष्ठ कौन्तेय युद्धाय कृतनिश्चयः॥ २.३७॥}
\translate{hatō vā prāpsyasi svargaṁ jitvā vā bhōkṣyasē mahīm |\\
tasmāduttiṣṭha kauntēya yuddhāya kṛtaniścayaḥ ‖ 2.37‖}
\cquote{If you are slain, you will attain heaven; if victorious, you will enjoy the earth. Therefore, arise, O son of Kunti (Arjuna), resolved to fight the battle.}
\slcol{\Index{सुखदुःखे समे कृत्वा} लाभालाभौ जयाजयौ।\\
ततो युद्धाय युज्यस्व नैवं पापमवाप्स्यसि॥ २.३८॥}
\translate{sukhaduḥkhē samē kṛtvā lābhālābhau jayājayau |\\
tatō yuddhāya yujyasva naivaṁ pāpamavāpsyasi ‖ 2.38‖}
\cquote{Making equal pleasure and pain, gain and loss, victory and defeat, then, resolved to fight the battle in this spirit, you will not incur sin.}
\slcol{\Index{एषा तेऽभिहिता} साङ्‍ख्ये बुद्धिर्योगे त्विमां शृणु।\\
बुद्ध्या युक्तो यया पार्थ कर्मबन्धं प्रहास्यसि॥ २.३९॥}
\translate{ēṣā tē'bhihitā sāṅ‍khyē buddhiryōgē tvimāṁ śṛṇu |\\
buddhyā yuktō yayā pārtha karmabandhaṁ prahāsyasi ‖ 2.39‖}
\cquote{This wisdom has been declared to you from the standpoint of discriminative knowledge (\emph{sankhya}).  Now hear it from the standpoint of  \emph{yoga} of wisdom, O Pārtha. Endowed with this understanding, you shall be freed from the bondage of action.}

\newpage
\begin{mananam}{\mananamfont {Mananam sloka - 38}}
\mananamtext What is my attitude towards any challenge I face in life? Am I equanimous towards any outcome or result? Despite putting in my best efforts, am I mentally averse to unpleasant outcomes and attached to favourable ones? Do I value learning to act in life without the motivation of seeking profits and avoiding losses, without wanting victory and rejecting defeat?
\end{mananam}
\WritingHand\enspace{\selfReflect\small{Self Reflection}}
\begin{inspiration}{\mananamfont Inspiration}
\mananamtext To seek to be victorious and to want to be happy is a natural inclination in all, as is the instinct to avoid the unpleasant and to never fail. Yet, for the truly free person, all actions are like drawing a line in water. He is mentally taintless and thus there is no \emph{karmic} bondage whatsoever.
\end{inspiration}
\newpage

\slcol{\Index{नेहाभिक्रमनाशोऽस्ति} प्रत्यवायो न विद्यते।\\
स्वल्पमप्यस्य धर्मस्य त्रायते महतो भयात्॥ २.४०॥}
\translate{nēhābhikramanāśō'sti pratyavāyō na vidyatē |\\
svalpamapyasya dharmasya trāyatē mahatō bhayāt ‖ 2.40‖}
\cquote{In this path (\emph{karma-yoga}), no effort is ever lost, nor is there any adverse effect. Even a little practice of this \emph{dharma} protects one from great fear.}
\slcol{\Index{व्यवसायात्मिका} बुद्धिरेकेह कुरुनन्दन।\\
बहुशाखा ह्यनन्ताश्च बुद्धयोऽव्यवसायिनाम्॥ २.४१॥}
\translate{vyavasāyātmikā buddhirēkēha kurunandana |\\
bahuśākhā hyanantāśca buddhayō'vyavasāyinām ‖ 2.41‖}
\cquote{For those who are resolute, the understanding is single-pointed, O son of Kuru (Arjuna). But the thoughts of the irresolute are many-branched and endless in their diversity.}
\slcol{\Index{यामिमां पुष्पितां} वाचं प्रवदन्त्यविपश्चितः।\\
वेदवादरताः पार्थ नान्यदस्तीति वादिनः॥ २.४२॥}
\translate{yāmimāṁ puṣpitāṁ vācaṁ pravadantyavipaścitaḥ |\\
vēdavādaratāḥ pārtha nānyadastīti vādinaḥ ‖ 2.42‖}
\cquote{For the undiscerning, delight lies in the flowery words of the vedas, O Partha(Arjuna).  Their speech being bound to rituals, declare that there is nothing beyond.}
\slcol{\Index{कामात्मानः स्वर्गपरा} जन्मकर्मफलप्रदाम्।\\
क्रियाविशेषबहुलां भोगैश्वर्यगतिं प्रति॥ २.४३॥}
\translate{kāmātmānaḥ svargaparā janmakarmaphalapradām |\\
kriyāviśēṣabahulāṁ bhōgaiśvaryagatiṁ prati ‖ 2.43‖}
\cquote{Driven by desire, intent on heaven as the goal, they proclaim ritual acts that yield rebirth and merit. Their words promise sensual enjoyment and dominion as rewards for various ritual activities.}
\newpage

\slcol{\Index{भोगैश्वर्यप्रसक्तानां} तयापहृतचेतसाम्।\\
व्यवसायात्मिका बुद्धिः समाधौ न विधीयते॥ २.४४॥}
\translate{bhōgaiśvaryaprasaktānāṁ tayāpahṛtacētasām |\\
vyavasāyātmikā buddhiḥ samādhau na vidhīyatē ‖ 2.44‖}
\cquote{For those whose minds are overpowered by attachment to enjoyment and dominion, resolute understanding does not arise. For such a one, steadiness in \emph{samadhi} cannot be established.}
\slcol{\Index{त्रैगुण्यविषया वेदा} निस्त्रैगुण्यो भवार्जुन।\\
निर्द्वन्द्वो नित्यसत्त्वस्थो निर्योगक्षेम आत्मवान्॥ २.४५॥}
\translate{traiguṇyaviṣayā vēdā nistraiguṇyō bhavārjuna |\\
nirdvandvō nityasattvasthō niryōgakṣēma ātmavān ‖ 2.45‖}
\cquote{The vedas deal with the three \emph{guṇas}. But O Arjuna, transcend the three \emph{guṇas}. Be free from dualities, ever firm in purity, indifferent to acquisition and preservation, and established in the Self.}
\slcol{\Index{यावानर्थ उदपाने} सर्वतःसम्प्लुतोदके।\\
तावान् सर्वेषु वेदेषु ब्राह्मणस्य विजानतः॥ २.४६॥}
\translate{yāvānartha udapānē sarvataḥsamplutōdakē |\\
tāvān sarvēṣu vēdēṣu brāhmaṇasya vijānataḥ ‖ 2.46‖}
\cquote{To the one who truly knows (Brahmana), all the vedas are of as much use as a small well is, when a vast flood surrounds on every side.}
\slcol{\Index{कर्मण्येवाधिकारस्ते} मा फलेषु कदाचन।\\
मा कर्मफलहेतुर्भूर्मा ते सङ्गोऽस्त्वकर्मणि॥ २.४७॥}
\translate{karmaṇyēvādhikārastē mā phalēṣu kadācana |\\
mā karmaphalahēturbhūrmā tē saṅgō'stvakarmaṇi ‖ 2.47‖}
\cquote{Your right is to action alone, never to its fruits. Let not the fruits of action be your motive, nor let your attachment be to inaction.}

%% \newgeometry{margin=0pt} % Apply margin only for this page
%% \thispagestyle{empty}
%% \begin{figure}
%% \centering
%% \includegraphics[width=\paperwidth, height=\paperheight, keepaspectratio]{../images/003.jpg}
%% \end{figure}
%% \restoregeometry % Restore original geometry settings
%%\newpage

\newpage
\begin{mananam}{\mananamfont {Mananam sloka - 45}}
\mananamtext Is my vision of life tainted by the pair of opposites, such as good and bad, beautiful and ugly, black and white, and so on? Am I ever judgmental towards others and myself? Do I cling to extreme views, which are a source of conflict within me and with others? Am I daring enough to rise above all views and judgments, not even holding to neutrality, but rather open and receptive to a dimension whose perspective is above all?
\end{mananam}
\WritingHand\enspace{\selfReflect\small{Self Reflection}}
\begin{inspiration}{\mananamfont Inspiration}
\mananamtext An immediate way to free yourself from all stress and anxieties is to take a few minutes in your daily schedule to rest the body and mind in a natural state of not doing, nor wanting to do, anything. It is neither meditating nor going to sleep. This state of being with your own Self is the most natural and deeply fulfilling state. It is free from the triads of \emph{sattva}, \emph{rajas} and \emph{tamas} which ever keep us outwardly engaged.
\end{inspiration}
\newpage

\slcol{\Index{योगस्थः कुरु} कर्माणि सङ्गं त्यक्त्वा धनञ्जय।\\
सिद्ध्यसिद्ध्योः समो भूत्वा समत्वं योग उच्यते॥ २.४८॥}
\translate{yōgasthaḥ kuru karmāṇi saṅgaṁ tyaktvā dhanañjaya |\\
siddhyasiddhyōḥ samō bhūtvā samatvaṁ yōga ucyatē ‖ 2.48‖}
\cquote{Established in \emph{yoga}, O Dhananjaya (Arjuna), perform your actions, abandoning attachment. Remain balanced in success and failure — such equanimity is called \emph{yoga}.}
\slcol{\Index{दूरेण ह्यवरं कर्म} बुद्धियोगाद्धनञ्जय।\\
बुद्धौ शरणमन्विच्छ कृपणाः फलहेतवः॥ २.४९॥}
\translate{dūrēṇa hyavaraṁ \emph{karma} buddhiyōgāddhanañjaya |\\
buddhau śaraṇamanviccha kṛpaṇāḥ phalahētavaḥ ‖ 2.49‖}
\cquote{Action driven by desire (\emph{sakama-karma}) is far inferior to action performed with the discipline of wisdom (of karma-yoga). Seek refuge in wisdom, for pitiable are those who act for reward.}
\slcol{\Index{बुद्धियुक्तो जहातीह} उभे सुकृतदुष्कृते।\\
तस्माद्योगाय युज्यस्व योगः कर्मसु कौशलम्॥ २.५०॥}
\translate{buddhiyuktō jahātīha ubhē sukṛtaduṣkṛtē |\\
tasmādyōgāya yujyasva yōgaḥ karmasu kauśalam ‖ 2.50‖}
\cquote{One established in wisdom-guided action, abandons both good and evil results (engages in \emph{nishkama-karma}). Therefore, devote yourself to \emph{yoga} — for \emph{yoga} is skill in action.}
\slcol{\Index{कर्मजं बुद्धियुक्ता} हि फलं त्यक्त्वा मनीषिणः।\\
जन्मबन्धविनिर्मुक्ताः पदं गच्छन्त्यनामयम्॥ २.५१॥}
\translate{karmajaṁ buddhiyuktā hi phalaṁ tyaktvā manīṣiṇaḥ |\\
janmabandhavinirmuktāḥ padaṁ gacchantyanāmayam ‖ 2.51‖}
\cquote{Those united with wisdom-guided action, abandoning the fruits of action,  are freed from the bondage of birth cycles. They attain the state beyond suffering, where supreme peace is found.}

\newpage
\begin{mananam}{\mananamfont {Mananam sloka - 47, 48}}
\mananamtext Can I even fathom the notion of doing something without seeking results and maintaining a state of equanimity, especially towards undesirable outcomes? Like most in society, am I constantly engaged in all interactions with others based on a transactional mindset? How can I train myself to give without taking? What are the little things in my daily life where I can begin this practice of “giving without taking” (before I feel comfortable in bringing this principle to bigger things in my life)?
\end{mananam}
\WritingHand\enspace{\selfReflect\small{Self Reflection}}
\begin{inspiration}{\mananamfont Inspiration}
\mananamtext Purifying the mind of all expectations is essential for attaining evenness of mind. Maintaining this evenness of mind in daily life is \emph{yoga}, and one who has achieved that state is a \emph{yogi}.
\end{inspiration}
\newpage

\newpage
\begin{mananam}{\mananamfont {Mananam sloka - 50}}
\mananamtext Am I equipped with the necessary skills to be successful not only in my outer life but in my inner life as well? Do I have the ability to handle emotions during stress, especially at work? Am I able to be harmonious in all my relationships? Am I able to be patient and persistent in my efforts towards my overall wellbeing? Do I have the Gita’s vision of going beyond worldly gains and losses to the state of mental equanimity?
\end{mananam}
\WritingHand\enspace{\selfReflect\small{Self Reflection}}
\begin{inspiration}{\mananamfont Inspiration}
\mananamtext The purpose of practicing \emph{yoga} is to equip us with necessary life skills. Just as we need driving skills to drive a car, so do we need life skills to function in life. Working skilfully, a \emph{karma-yogi} becomes free from the bondage which would be otherwise induced by actions.
\end{inspiration}
\newpage


\slcol{\Index{यदा ते मोहकलिलं} बुद्धिर्व्यतितरिष्यति।\\
तदा गन्तासि निर्वेदं श्रोतव्यस्य श्रुतस्य च॥ २.५२॥}
\translate{yadā tē mōhakalilaṁ buddhirvyatitariṣyati |\\
tadā gantāsi nirvēdaṁ śrōtavyasya śrutasya ca ‖ 2.52‖}
\cquote{When your wisdom crosses beyond the mire of delusion, you will become indifferent to what has been heard and what is yet to be heard.}
\slcol{\Index{श्रुतिविप्रतिपन्ना ते} यदा स्थास्यति निश्चला।\\
समाधावचला बुद्धिस्तदा योगमवाप्स्यसि॥ २.५३॥}
\translate{śrutivipratipannā tē yadā sthāsyati niścalā |\\
samādhāvacalā buddhistadā yōgamavāpsyasi ‖ 2.53‖}
\cquote{When your wisdom, no longer bewildered by conflicting teachings, stands firm and unmoving, fixed in Self (\emph{samadhi}), then you will attain \emph{yoga} (the state beyond doubt).}
\slcol{अर्जुन उवाच —\\
\Index{स्थितप्रज्ञस्य का भाषा} समाधिस्थस्य केशव।\\
स्थितधीः किं पृभाषेत किमासीत व्रजेत किम्॥ २.५४॥}
\translate{arjuna uvāca —\\
sthitaprajñasya kā bhāṣā samādhisthasya kēśava |\\
sthitadhīḥ kiṁ pṛbhāṣēta kimāsīta vrajēta kim ‖ 2.54‖}
\cquote{Arjuna said -\\
O Keshava (Krishna), what are the characteristics of one whose wisdom is steady (\emph{sthita-prajña}), and absorbed in \emph{samadhi}? How does such a person speak, sit, and walk?}


\newpage
\begin{mananam}{\mananamfont {Mananam sloka - 52, 53}}
\mananamtext Has my mind attained a degree of steadiness on the path of \emph{yoga}? Am I still unclear and confused on the path? Do I have a means to clarify my doubts? Which is the primary source—teaching or teacher—in which I seek refuge? Am I feeling greater attunement with that inner principle of the \emph{Guru} through these external sources of refuge?
\end{mananam}
\WritingHand\enspace{\selfReflect\small{Self Reflection}}
\begin{inspiration}{\mananamfont Inspiration}
\mananamtext The purpose of the teacher and the scriptures is to shake us out of delusion. Confusion and challenges are a part of the spiritual journey as we let go of our worldly thinking patterns and seek refuge in the higher reality. As we thus progress, we begin to gain clarity on the path.
\end{inspiration}
\newpage

\begin{mananam}{\mananamfont {Mananam sloka - 54}}
\mananamtext What do I understand to be the state of realization? Do I know about the different degrees of realization of the saints? What is the next stage in my own spiritual evolution? What stage can I sincerely aspire toward? Is my understanding of ultimate liberation and freedom motivating me to make mental and spiritual progress in my life?
\end{mananam}
\WritingHand\enspace{\selfReflect\small{Self Reflection}}
\begin{inspiration}{\mananamfont Inspiration}
\mananamtext In every field of life, we are motivated by people who are successful. Similarly, the saints and sages are those who have attained a state that is hard for many of us to fathom, as it is not just an outer state, but an inner state. Gaining a proper understanding of their inner state helps us to refrain from erroneous judgements and conclusions.
\end{inspiration}
\newpage


\slcol{श्रीभगवानुवाच —\\
\Index{प्रजहाति यदा} कामान्सर्वान्पार्थ मनोगतान्।\\
आत्मन्येवात्मना तुष्टः स्थितप्रज्ञस्तदोच्यते॥ २.५५॥}
\translate{śrībhagavānuvāca —\\
prajahāti yadā kāmānsarvānpārtha manōgatān |\\
ātmanyēvātmanā tuṣṭaḥ sthitaprajñastadōcyatē ‖ 2.55‖}
\cquote{Sri Bhagavan said: When one completely abandons all desires that arise in the mind, and when the Self is content in the Self alone, then that person is said to be of steady wisdom.}
\slcol{\Index{दुःखेष्वनुद्विग्नमनाः} सुखेषु विगतस्पृहः।\\
वीतरागभयक्रोधः स्थितधीर्मुनिरुच्यते॥ २.५६॥}
\translate{duḥkhēṣvanudvignamanāḥ sukhēṣu vigataspṛhaḥ |\\
vītarāgabhayakrōdhaḥ sthitadhīrmunirucyatē ‖ 2.56‖}
\cquote{One whose mind is not shaken by adversity, who does not crave pleasures, and who is free from attachment, fear, and anger; such a person is said to have a steady intellect.}
\slcol{\Index{यः सर्वत्रानभिस्नेहस्त}त्तत्प्राप्य शुभाशुभम्।\\
नाभिनन्दति न द्वेष्टि तस्य प्रज्ञा प्रतिष्ठिता॥ २.५७॥}
\translate{yaḥ sarvatrānabhisnēhastattatprāpya śubhāśubham |\\
nābhinandati na dvēṣṭi tasya prajñā pratiṣṭhitā ‖ 2.57‖}
\cquote{One who is without attachment, who neither rejoices when good comes nor hates when evil comes—his wisdom is firmly established.}

\newpage
\begin{mananam}{\mananamfont {Mananam sloka - 55}}
\mananamtext What are my lower desires in life that keep me from making spiritual progress? How can I renounce the mental habits and passions that pull me downward? What does it mean to find happiness and satisfaction in one’s own Self?
\end{mananam}
\WritingHand\enspace{\selfReflect\small{Self Reflection}}
\begin{inspiration}{\mananamfont Inspiration}
\mananamtext A true \emph{yogi} or \emph{sannyasi} is one who is not dependent on someone or something to be happy. There are no worldly desires to hanker after as he or she finds all of them fulfilled within their own Self.
\end{inspiration}
\newpage


\newpage
\begin{mananam}{\mananamfont {Mananam sloka - 56, 57}}
\mananamtext Do failures and losses bring me down mentally? Do I take it personally? Do I allow it to affect my self-confidence? How about successes and gains? Do I get overly excited? Does it make me arrogant and condescending towards others? Do I constantly judge life circumstances as being either good or bad?
\end{mananam}
\WritingHand\enspace{\selfReflect\small{Self Reflection}}
\begin{inspiration}{\mananamfont Inspiration}
\mananamtext For common people, the outer circumstances of life tend to sway their mental states accordingly. Thus arise the natural inclination towards positive outcomes in life, and aversion towards negative. For a \emph{yogi}, the goal of all spiritual practices is not to gain favourable life circumstances, but rather, that unshakeable steadiness in mind.
\end{inspiration}
\newpage

\slcol{\Index{यदा संहरते चायं} कूर्मोऽङ्गानीव सर्वशः।\\
इन्द्रियाणीन्द्रियार्थेभ्यस्तस्य प्रज्ञा प्रतिष्ठिता॥ २.५८॥}
\translate{yadā saṁharatē cāyaṁ kūrmō'ṅgānīva sarvaśaḥ |\\
indriyāṇīndriyārthēbhyastasya prajñā pratiṣṭhitā ‖ 2.58‖}
\cquote{When a person completely withdraws his senses from the sense objects, just as a tortoise withdraws its limbs from all sides into its shell,  his wisdom is firmly established.}
\slcol{\Index{विषया विनिवर्तन्ते} निराहारस्य देहिनः।\\
रसवर्जं रसोऽप्यस्य परं दृष्ट्वा निवर्तते॥ २.५९॥}
\translate{viṣayā vinivartantē nirāhārasya dēhinaḥ |\\
rasavarjaṁ rasō'pyasya paraṁ dṛṣṭvā nivartatē ‖ 2.59‖}
\cquote{Objects of senses turn away from one who has abstained from them, leaving behind the taste for them; upon realising the Supreme, even the taste falls away.}
\slcol{\Index{यततो ह्यपि कौन्तेय} पुरुषस्य विपश्चितः।\\
इन्द्रियाणि प्रमाथीनि हरन्ति प्रसभं मनः॥ २.६०॥}
\translate{yatatō hyapi kauntēya puruṣasya vipaścitaḥ |\\
indriyāṇi pramāthīni haranti prasabhaṁ manaḥ ‖ 2.60‖}
\cquote{The turbulent sense organs, O Kaunteya (Arjuna, son of Kunti), violently snatch away the mind of even a wise person, though he may be earnestly striving to control them.}
\slcol{\Index{तानि सर्वाणि संयम्य} युक्त आसीत मत्परः।\\
वशे हि यस्येन्द्रियाणि तस्य प्रज्ञा प्रतिष्ठिता॥ २.६१॥}
\translate{tāni sarvāṇi saṁyamya yukta āsīta matparaḥ |\\
vaśē hi yasyēndriyāṇi tasya prajñā pratiṣṭhitā ‖ 2.61‖}
\cquote{The one who has restrained all the senses should sit steadfast, intent on Me as the supreme; his wisdom is firmly established, whose senses are under control.}

\newpage
\begin{mananam}{\mananamfont {Mananam sloka - 58}}
\mananamtext How much sense of control do I have? Do I tend to overindulge in sensual pleasures? Am I aware of the external triggers that draw my mind towards the senses? Am I aware of the internal triggers of thoughts and memories that present the objects of attraction? Have I attempted any methods of withdrawing the mind from the senses at will? If not, how may I develop it?
\end{mananam}
\WritingHand\enspace{\selfReflect\small{Self Reflection}}
\begin{inspiration}{\mananamfont Inspiration}
\mananamtext It takes conscious training to withdraw the naturally-outgoing senses inward. Mere wishing and knowledge of the harm of addictions is not adequate. Taking consistent and incremental steps is the best approach to cultivate new positive habits.
\end{inspiration}
\newpage


\slcol{\Index{ध्यायतो विषयान्पुंसः} सङ्गस्तेषूपजायते। \\
सङ्गात्सञ्जायते कामः कामात्क्रोधोऽभिजायते॥ २.६२॥}
\translate{dhyāyatō viṣayānpuṁsaḥ saṅgastēṣūpajāyatē |\\
saṅgātsañjāyatē kāmaḥ kāmātkrōdhō'bhijāyatē ‖ 2.62‖}
\cquote{When a person dwells on sense objects, attachment to them arises. From attachment, desire is born; and when desire remains unfulfilled, it gives rise to anger.}
\slcol{\Index{क्रोधाद्भवति संमोहः} संमोहात्स्मृतिविभ्रमः।\\
स्मृतिभ्रंशाद्बुद्धिनाशो बुद्धिनाशात्प्रणश्यति॥ २.६३॥}
\translate{krōdhādbhavati saṁmōhaḥ saṁmōhātsmṛtivibhramaḥ |\\
smṛtibhraṁśādbuddhināśō buddhināśātpraṇaśyati ‖ 2.63‖}
\cquote{Anger leads to delusion; and delusion results in the loss of memory. From loss of memory, the intellect is destroyed – and with the destruction of intellect, one falls into complete ruin.}
\slcol{\Index{रागद्वेषवियुक्तैस्तु} विषयानिन्द्रियैश्चरन्।\\
आत्मवश्यैर्विधेयात्मा प्रसादमधिगच्छति॥ २.६४॥}
\translate{rāgadvēṣaviyuktaistu viṣayānindriyaiścaran |\\
ātmavaśyairvidhēyātmā prasādamadhigacchati ‖ 2.64‖}
\cquote{But a self-controlled person, free from attachment and aversion, who moves among sense objects with controlled senses and a disciplined mind, attains tranquillity.}
\slcol{\Index{प्रसादे सर्वदुःखानां} हानिरस्योपजायते।\\
प्रसन्नचेतसो ह्याशु बुद्धिः पर्यवतिष्ठते॥ २.६५॥}
\translate{prasādē sarvaduḥkhānāṁ hānirasyōpajāyatē |\\
prasannacētasō hyāśu buddhiḥ paryavatiṣṭhatē ‖ 2.65‖}
\cquote{Through tranquillity, all suffering is eradicated. For one with a serene mind, the intellect becomes steady.}


\newpage
\begin{center}
  {\Large Ladder of Fall} \\(Sloka 62 and 63)
\end{center}
\tikzstyle{startstop} = [rectangle, rounded corners, 
minimum width=4cm, 
minimum height=1cm,
text centered, 
draw=black,
align=center,
fill=red!30]

\tikzstyle{io} = [trapezium, 
trapezium stretches=true, % A later addition
trapezium left angle=70, 
trapezium right angle=110, 
minimum width=3cm, 
minimum height=1cm, text centered, 
draw=black, fill=blue!30]

\tikzstyle{process} = [rectangle, 
minimum width=4cm, 
minimum height=1cm, 
text centered, 
text width=3cm, 
draw=black, 
fill=orange!30]

\tikzstyle{etbox} = [rectangle, 
minimum width=4cm, 
minimum height=1cm, 
text centered, 
text width=3cm, 
draw=white, 
fill=none]

\tikzstyle{decision} = [diamond, 
minimum width=3cm, 
minimum height=1cm, 
text centered, 
draw=black, 
fill=green!30]
\tikzstyle{arrow} = [thick,->,>=stealth]

\begin{tikzpicture}[node distance=2cm]
\node (start) [startstop] {Thinking about Objects\\ध्यायतो विशयान्पुम्सः}; \node(l1)[etbox, left of =start,xshift=-2.5cm]{Levels of Thinking};
\node(pro1) [process, below of =start] {Mentally bound to objects}; \node(l2)[etbox,left of =pro1,xshift=-2.5cm]{Difficult to control senses};
\node(pro2) [process, below of =pro1] {Greed due to mental binding}; \node(l3)[etbox, left of =pro2,xshift=-2.5cm]{Mental Agony};
\node(pro3) [process, below of =pro2] {Anger due to unfulfilled desires}; \node(l4)[etbox, left of =pro3,xshift=-2.5cm]{Control is Impossible};
\node(pro4) [process, below of =pro3] {Anger leading to mental disarray}; \node(l5)[etbox, left of =pro4,xshift=-2.5cm]{Constant Fall};
\node(pro5) [process, below of =pro4] {Forgetfullness due to mental disarray};
\node(pro6) [process, below of =pro5] {Destroyed mind due to Forgetfullness};
\node (stop) [startstop, below of=pro6] {End of Life};

\draw [arrow] (start) -- (pro1);
\draw [arrow] (pro1) -- (pro2);
\draw [arrow] (pro2) -- (pro3);
\draw [arrow] (pro3) -- (pro4);
\draw [arrow] (pro4) -- (pro5);
\draw [arrow] (pro5) -- (pro6);
\draw [arrow] (pro6) -- (stop);

\end{tikzpicture}
\newpage
\begin{mananam}{\mananamfont {Mananam sloka - 62, 63}}
\mananamtext Am I aware of the sequence of this ‘ladder of falls’ in my daily life? Can I relate to how brooding on objects or people has led to attachment? Can I relate to events of anger arising when desires were obstructed?
Looking back to events when I was angry, can I see how they led to a delusional state and loss of the memory of my good intentions? And, from that, the loss of ability to discern good from bad, right from wrong, etc? Can I also learn by observing others’ falling to this sequence in their lives?
\end{mananam}
\WritingHand\enspace{\selfReflect\small{Self Reflection}}
\begin{inspiration}{\mananamfont Inspiration}
\mananamtext A clear sequence of the “ladder of falls” from the Gita and how one thing leads to another is a very valuable insight. All around us we can see examples of people falling victim to this chain of events. The more progressively lower on the ladder our fall occurs, the harder it is to reverse it.
\end{inspiration}
\newpage

\slcol{\Index{नास्ति बुद्धिरयुक्तस्य} न चायुक्तस्य भावना।\\
न चाभावयतः शान्तिरशान्तस्य कुतः सुखम्॥ २.६६॥}
\translate{nāsti buddhirayuktasya na cāyuktasya bhāvanā |\\
na cābhāvayataḥ śāntiraśāntasya kutaḥ sukham ‖ 2.66‖}
\cquote{For the unsteady, there is no wisdom; and for such a one, no contemplation is possible. Without contemplation, there is no peace and without peace how can there be happiness?}
\slcol{\Index{इन्द्रियाणां हि} चरतां यन्मनो'नुविधीयते।\\
तदस्य हरति प्रज्ञां वायुर्नावमिवाम्भसि॥ २.६७॥}
\translate{indriyāṇāṁ hi caratāṁ yanmanō'nuvidhīyatē |\\
tadasya harati prajñāṁ vāyurnāvamivāmbhasi ‖ 2.67‖}
\cquote{Just as a boat on the water is swept away by a strong wind, when the mind yields to the wandering senses, one’s wisdom gets carried away.}
\slcol{\Index{तस्माद्यस्य महाबाहो} निगृहीतानि सर्वशः।\\
इन्द्रियाणीन्द्रियार्थेभ्यस्तस्य प्रज्ञा प्रतिष्ठिता॥ २.६८॥}
\translate{tasmādyasya mahābāhō nigṛhītāni sarvaśaḥ |\\
indriyāṇīndriyārthēbhyastasya prajñā pratiṣṭhitā ‖ 2.68‖}
\cquote{Therefore, O mighty-armed Arjuna, one whose senses are completely withdrawn from the sense objects – his wisdom is firmly established.}
\slcol{\Index{या निशा सर्वभूतानां} तस्यां जागर्ति संयमी।\\
यस्यां जाग्रति भूतानि सा निशा पश्यतो मुनेः॥ २.६९॥}
\translate{yā niśā sarvabhūtānāṁ tasyāṁ jāgarti saṁyamī |\\
yasyāṁ jāgrati bhūtāni sā niśā paśyatō munēḥ ‖ 2.69‖}
\cquote{That which is night for all beings is the awakening for the self-controlled; and that which is awakening for all beings is the night for the sage who sees the truth.}

\newpage
\begin{mananam}{\mananamfont {Mananam sloka - 67, 68}}
\mananamtext How good is my ability to discern between that which is good for me in the long term versus that which is good only in the short term? Am I able to discern habits that help me in my life and consciously work towards strengthening them?
In any indulgence, there are two processes – that of the mind dwelling on objects (which requires \emph{shama}, the mind’s restraint) and that of the physical senses grasping at them (which requires \emph{dhama}, the senses’ restraint). Since these happen so fast, most are not able to separate them. Do you recognize in your own life certain situations where you can see these two processes at work? Which of these do you need to work on the most?
\end{mananam}
\WritingHand\enspace{\selfReflect\small{Self Reflection}}
\begin{inspiration}{\mananamfont Inspiration}
\mananamtext There are many in this world who are a slave to their senses! The modern world not only considers it natural but even glorifies such a life. But all those who revert to such a hedonistic life are sure to become miserable in the long run. Restraint of the senses is what most religions admonish for the novices. For those deeply committed to progress, mental restraint is the higher practice.
\end{inspiration}
\newpage


\newpage
\begin{mananam}{\mananamfont {Mananam sloka - 69}}
\mananamtext Do I look forward to sleep as a means of escape from daily life? Do I have the sense of purpose in life that wakes me up every day with enthusiasm? Do the tendencies of laziness and boredom dominate my daily life?
Do I consider my waking state to be the ultimate reality, or can I perceive anything higher than that? Am I caught in the rat race of a life influenced by societal patterns? How can I cultivate the sage-like attitude of heightened awareness, in both the waking and sleep states?
\end{mananam}
\WritingHand\enspace{\selfReflect\small{Self Reflection}}
\begin{inspiration}{\mananamfont Inspiration}
\mananamtext Many of us tend to live an unconscious zombie-like life, dragged down by our sense-mind and impulses. Such living is as good as being asleep. Some others are driven by their desires and passions, but lack the discernment as to wherein lies their ultimate happiness. To be truly awake is to learn the art of conscious living.
\end{inspiration}
\newpage

\slcol{\Index{आपूर्यमाणमचलप्रतिष्ठं} समुद्रमापः प्रविशन्ति यद्वत्। \\
तद्वत्कामा यं प्रविशन्ति सर्वे स शान्तिमाप्नोति न कामकामी॥ २.७०॥}
\translate{āpūryamāṇamacalapratiṣṭhaṁ samudramāpaḥ\\praviśanti yadvat | \\
tadvatkāmā yaṁ praviśanti sarvē sa śāntimāpnōti\\na kāmakāmī ‖ 2.70‖}
\cquote{He attains peace into whom all desires enter as waters enter the ocean, which filled from all sides, remains unmoved; but not the one who hankers after the desires.}
\slcol{\Index{विहाय कामान्यः} सर्वान्पुमांश्चरति निःस्पृहः।\\
निर्ममो निरहङ्कारः स शान्तिमधिगच्छति॥ २.७१॥}
\translate{vihāya kāmānyaḥ sarvānpumāṁścarati niḥspṛhaḥ | \\
nirmamō nirahaṅkāraḥ sa śāntimadhigacchati ‖ 2.71‖}
\cquote{He who renounces all desires, moves about without longing, and is free from the notions of ‘me’ and ‘mine’ – such a one attains peace.}
\slcol{\Index{एषा ब्राह्मी स्थितिः} पार्थ नैनां प्राप्य विमुह्यति।\\
स्थित्वास्यामन्तकालेऽपि ब्रह्मनिर्वाणमृच्छति॥ २.७२॥}
\translate{ēṣā brāhmī sthitiḥ pārtha naināṁ prāpya vimuhyati |\\
sthitvāsyāmantakālē'pi brahmanirvāṇamṛcchati ‖ 2.72‖}
\cquote{O Partha, this is the state of being established in Brahman – attaining which, one is never deluded. Abiding in it even at the end of life, one attains liberation in Brahman.}

\newpage
\begin{mananam}{\mananamfont {Mananam sloka - 70, 71}}
\mananamtext How do I relate to goals and desires in my life? Can I work to fulfil my noble goals, but remain unshaken by the final results? What is the nature of my desires – selfish, harmful, egoistic?
In seeking goals, am I really seeking for acceptance, self-esteem, self-acknowledgement etc.? How can I learn in my life situation to turn the attitude of seeking the rewards of outward possessions to that of seeking inward gratification? How can I rephrase my goal statement to be universally beneficial?
\end{mananam}
\WritingHand\enspace{\selfReflect\small{Self Reflection}}
\begin{inspiration}{\mananamfont Inspiration}
\mananamtext The various needs and demands of life are imposing on everyone. But the wiser ones don’t focus their actions for the sake of sensual gratification or material possession, as they know that such longings can never come to an end.
Only those who realize their being to be of infinite nature know true contentment.
\end{inspiration}

\newpage
\begin{center}
{\devfont ॐ तत्सदिति श्रीमद्भगवद्गीतासूपनिषत्सु ब्रह्मविद्यायां योगशास्त्रे\\
श्रीकृष्णार्जुनसंवादे साङ्ख्ययोगो नाम द्वितीयोऽध्यायः॥ \\}
\chapEndSlokaTranslate{ōṃ tatsaditi śrīmadbhagavadgītās ūpaniṣatsu \\brahmavidyāyāṃ yōgaśāstrē\\ śrīkṛṣṇārjunasaṃvādē sāṅkhya-\emph{yoga}\\nāma  dvitīyo'dhyāyaḥ‖}
\end{center}